\documentclass[11pt,a4paper]{article}
\usepackage[margin=2.4cm]{geometry}
\usepackage{amsmath,amsthm,mathtools}
\usepackage{fontspec}
\usepackage{unicode-math}
\setmainfont{STIX2Text}[Path=fonts/, Extension=.otf,
  UprightFont=*-Regular, ItalicFont=*-Italic,
  BoldFont=*-Bold, BoldItalicFont=*-BoldItalic]
\setmathfont{STIX2Math}[Path=fonts/, Extension=.otf]
\usepackage[protrusion=true,expansion=false]{microtype}
\usepackage{booktabs,enumitem}
\usepackage{placeins}
\usepackage{tikz}
\usetikzlibrary{arrows.meta,positioning,calc,fit,backgrounds}
\definecolor{tierspend}{RGB}{200,60,60}
\definecolor{tierfull}{RGB}{210,140,40}
\definecolor{tierin}{RGB}{60,130,180}
\definecolor{tierrecv}{RGB}{70,150,90}
\usepackage[psdextra,hidelinks]{hyperref}

\emergencystretch=3em
\hbadness=10000
\hfuzz=2pt

\newtheorem{theorem}{Theorem}[section]
\newtheorem{lemma}[theorem]{Lemma}
\newtheorem{proposition}[theorem]{Proposition}
\newtheorem{corollary}[theorem]{Corollary}
\theoremstyle{definition}
\newtheorem{definition}[theorem]{Definition}
\newtheorem{construction}[theorem]{Construction}
\newtheorem{assumption}[theorem]{Assumption}
\newtheorem{remark}[theorem]{Remark}
\newtheorem{example}[theorem]{Example}

% kind-coded theorem blocks, palette shared with the figures and the website:
% definitions/constructions blue (tierin), assumptions amber (tierfull),
% proved statements a neutral bar; remarks and examples run unframed
\usepackage{mdframed}
\providecolor{tierin}{RGB}{60,130,180}
\providecolor{tierfull}{RGB}{210,140,40}
\mdfdefinestyle{kindbar}{topline=false,bottomline=false,rightline=false,
  linewidth=1.5pt,innerleftmargin=9pt,innerrightmargin=4pt,
  innertopmargin=5pt,innerbottommargin=6pt,leftmargin=0pt,rightmargin=0pt,
  skipabove=10pt,skipbelow=10pt}
\surroundwithmdframed[style=kindbar,linecolor=tierin!70,backgroundcolor=tierin!4]{definition}
\surroundwithmdframed[style=kindbar,linecolor=tierin!70,backgroundcolor=tierin!4]{construction}
\surroundwithmdframed[style=kindbar,linecolor=tierfull!80,backgroundcolor=tierfull!5]{assumption}
\surroundwithmdframed[style=kindbar,linecolor=black!30]{theorem}
\surroundwithmdframed[style=kindbar,linecolor=black!30]{lemma}
\surroundwithmdframed[style=kindbar,linecolor=black!30]{proposition}
\surroundwithmdframed[style=kindbar,linecolor=black!30]{corollary}

\usepackage{fancyhdr}
\pagestyle{fancy}
\fancyhf{}
\fancyhead[L]{\footnotesize\scshape The Zcash Arboretum}
\fancyhead[R]{\footnotesize\itshape\nouppercase{\leftmark}}
\fancyfoot[C]{\footnotesize\thepage}
\renewcommand{\headrulewidth}{0.3pt}
\renewcommand{\sectionmark}[1]{\markboth{\thesection\ \ #1}{}}

\newcommand{\F}{\mathbb{F}}
\newcommand{\Z}{\mathbb{Z}}
\newcommand{\N}{\mathbb{N}}
\newcommand{\G}{\mathbb{G}}
\newcommand{\E}{\mathbb{E}}
\newcommand{\PP}{\mathbb{P}}
\renewcommand{\Pr}{\mathrm{Pr}}
\newcommand{\Fp}{\F_p}
\newcommand{\Fq}{\F_q}
\newcommand{\ip}[2]{\langle #1,\, #2 \rangle}
\newcommand{\Comm}{\mathsf{Commit}}
\newcommand{\Open}{\mathsf{Open}}
\newcommand{\Setup}{\mathsf{Setup}}
\newcommand{\Verify}{\mathsf{Verify}}
\newcommand{\negl}{\mathsf{negl}}
\newcommand{\poly}{\mathsf{poly}}
\newcommand{\PPT}{\mathsf{PPT}}
\newcommand{\Adv}{\mathcal{A}}
\newcommand{\Prov}{\mathcal{P}}
\newcommand{\Ver}{\mathcal{V}}
\newcommand{\Sim}{\mathcal{S}}
\newcommand{\Ext}{\mathcal{E}}
\newcommand{\Rel}{\mathcal{R}}
\newcommand{\Lang}{\mathcal{L}}
\newcommand{\nfsym}{\mathsf{nf}}
\newcommand{\cmsym}{\mathsf{cm}}
\newcommand{\cmx}{\mathsf{cmx}}
\newcommand{\cvsym}{\mathsf{cv}}
\newcommand{\rk}{\mathsf{rk}}
\newcommand{\Pallas}{\ensuremath{\mathsf{Pallas}}}
\newcommand{\Vesta}{\ensuremath{\mathsf{Vesta}}}
\newcommand{\Ep}{\ensuremath{\mathcal{E}}}
\newcommand{\NoteCommit}{\ensuremath{\mathsf{NoteCommit}}}
\newcommand{\Sinsemilla}{\ensuremath{\mathsf{Sinsemilla}}}
\newcommand{\ShortCommit}{\ensuremath{\mathsf{ShortCommit}}}
\newcommand{\CommitIvk}{\ensuremath{\mathsf{Commit}^{\mathsf{ivk}}}}
\newcommand{\Extract}{\ensuremath{\mathsf{Extract}}}
\newcommand{\Acc}{\ensuremath{\mathsf{Acc}}}
\newcommand{\GroupHash}{\ensuremath{\mathsf{GroupHash}}}
\newcommand{\ask}{\ensuremath{\mathsf{ask}}}
\newcommand{\aksym}{\ensuremath{\mathsf{ak}}}
\newcommand{\nksym}{\ensuremath{\mathsf{nk}}}
\newcommand{\ivk}{\ensuremath{\mathsf{ivk}}}
\newcommand{\fvk}{\ensuremath{\mathsf{fvk}}}
\newcommand{\ovk}{\ensuremath{\mathsf{ovk}}}
\newcommand{\dk}{\ensuremath{\mathsf{dk}}}
\newcommand{\pkd}{\ensuremath{\mathsf{pk_d}}}
\newcommand{\gd}{\ensuremath{\mathsf{g_d}}}
\newcommand{\rcv}{\ensuremath{\mathsf{rcv}}}
\newcommand{\rcm}{\ensuremath{\mathsf{rcm}}}
\newcommand{\rivk}{\ensuremath{\mathsf{r_{ivk}}}}
\newcommand{\rsk}{\ensuremath{\mathsf{rsk}}}

\renewenvironment{abstract}{%
  \small\begin{center}{\bfseries\abstractname}\end{center}\medskip\noindent\ignorespaces
}{\par\medskip}


\title{\textbf{\Huge ZSA Guide}\\[6pt]\large Issuance, private transfer, and public supply}
\author{m@rek.onl}
\date{}
\begin{document}
\maketitle
\thispagestyle{empty}

\begin{abstract}
A shielded asset is not a number with a hidden text label. Its identity
must survive issuance, commitments, transfers, splitting, and burning,
without letting an adversary invent algebraic exchange rates between
unrelated assets. This volume constructs that chain of obligations.
It derives the issuer's signature, the public issuance state machine,
the private per-asset balance equation, and the membership rule that
makes zero-value split inputs safe. The endpoint is the mathematical
and validation design in one pinned pair of draft ZIPs. Where those
drafts leave their transaction carrier or note-format instantiation
unfinished, the boundary remains visible.
\end{abstract}

\clearpage
\tableofcontents
\clearpage

\section{The asset ledger contract}
\label{sec:zsa-contract}

The integer, field, and point-encoding language used here comes from
\emph{Math Guide}, \S\emph{Prime fields and modular arithmetic} and
\S\emph{Pasta curves and their encodings}. Probability, including
uniform randomness, is constructed in \emph{Math Guide},
\S\emph{Probability and computation}. \emph{Crypto Guide},
\S\emph{Security contracts and reductions}, defines computational
security and failure probabilities. Its
\S\emph{Groups, commitments, and Sinsemilla} defines hiding and binding
and constructs the commitments used by the inherited note relation;
\S\emph{Knowledge, simulation, and Schnorr} defines knowledge and
extraction.

\subsection{What changes when there is more than one asset}

In the native payment protocol, every amount uses the same unit:
a zatoshi. If one input holds \(80\,000\), outputs hold \(50\,000\)
and \(20\,000\), and the fee is \(10\,000\), the equation
\[
  80\,000=50\,000+20\,000+10\,000
\]
has a meaning before any cryptography is applied.

For two assets \(A\) and \(B\), the expression \(80A=50B+30A\)
is not a balanced payment. No exchange rate is part of the shielded
ledger's addition operation. A valid transaction must conserve each
asset separately, except for explicitly authorised creation or public
destruction of that same asset.

The design separates three operations.
\begin{enumerate}
  \item \emph{Issuance} creates custom-asset notes. The asset,
    recipient, amount, and note randomness are publicly available,
    and an issuer signature authorises the creation.
  \item \emph{Transfer} consumes notes and creates notes. The proof
    hides their assets and amounts while enforcing per-asset balance.
  \item \emph{Burn} removes a public amount of a named custom asset
    from the shielded ledger. The proof and binding signature account
    for that removal.
\end{enumerate}
An issuer controls further creation of its assets, not the spend
authorisation keys of holders. The native asset retains its existing
value base and pays transaction fees. A custom asset cannot be moved
into a transparent native-value field by relabelling its units.

\subsection{The selected endpoint}

The rule set here is
\href{https://github.com/zcash/zips/blob/69610984109078075e7e64989ecf89d63a259b97/zips/zip-0226.rst}
{ZIP 226} and
\href{https://github.com/zcash/zips/blob/69610984109078075e7e64989ecf89d63a259b97/zips/zip-0227.rst}
{ZIP 227} at commit \texttt{69610984}.
They are draft specifications, not evidence that these rules are active
on a network. Their Orchard definitions incorporate the ZIP 2005
changes at that same source pin.

The asset mechanisms are substantially specified, but the pair still
references the withdrawn ZIP 230 transaction carrier and ZIP 246
digest scheme. Its ZSA note lead byte is a placeholder. Consequently,
this volume does not claim to define an interoperable, deployable
ZSA transaction encoding. It constructs the specified state and proof
contracts and identifies exactly which missing interfaces prevent
turning them into a complete byte-level validator. The version-six
transaction in the core volumes is not silently extended with these
fields.

\subsection{Dependencies and notation}

\emph{Ironwood Guide}, \S\emph{Actions and the proof relation}, supplies the
one-input/one-output Action, the note commitment, Merkle membership,
nullifier, and ownership relation. Its
\S\emph{Balance, digests, and authorisation} supplies the binding
signature and the distinction between transaction effects and
authorising data. We change those constructions only where stated.

\emph{Math Guide}, \S\emph{Pasta curves and their encodings}, supplies Pallas.
Write \(\mathbb P\) for its prime-order group, \(p\) for its base-field
modulus, and \(q\) for its group order. Thus note values such as
\(\rho,\psi,\nfsym\) lie in \(\F_p\), whereas commitment randomness
and signature scalars lie in \(\F_q\). These two fields must not be
interchanged.

Write \(V\) for the existing Orchard native-value base and \(R_v\)
for its value-commitment randomness base. An asset base \(A\) is a
nonidentity Pallas point. The point's canonical encoding \(A^*\)
is 32 bytes; when fed to a bit-oriented commitment, these same bytes
are unpacked least-significant bit first. \(\operatorname{enc}_k(x)\)
means exactly \(k\) bytes of little-endian integer encoding.
The shared notation \(\mathrm{LE}_k(x)\) instead denotes exactly
\(k\) bits, least-significant bit first.

\subsection{State is branch-local}

A full validator keeps the existing note tree, recognised anchors,
and spent-nullifier set. It additionally keeps an issuance map
\[
  I[A]=(\mathsf{balance},\mathsf{final},\mathsf{reference}).
\]
The map is part of the state of a particular chain branch. Processing
a transaction produces a candidate next state; any failed check
discards the entire candidate. A reorganisation restores the state
at the common ancestor and replays the replacement branch.

In particular, an issuance finalisation is irreversible \emph{along
a fixed accepted branch}, not immune to a reorganisation of that
branch. No extra finality assumption is hidden in the word
\(\mathsf{final}\).

\section{Asset identity and issuance authority}
\label{sec:zsa-identity}

\subsection{Deriving a separate issuance key}

The issuance secret is independent of the protocol role of a note
spending key. A holder may authorise a payment without gaining
authority to issue more units, and an issuer may sign issuance
transactions without learning holders' spend keys.

The derivation reuses the hardened-only machinery in \emph{Wallet Guide},
\S\emph{From seed to account}, with a different master domain and
path. Let the seed \(S\) contain between 32 and 252 bytes; it should
contain at least 256 bits of entropy. For the guessing threat here,
a conservative meaning is that no one seed occurs with probability
greater than \(2^{-256}\), called at least 256 bits of
\emph{min-entropy}. Merely storing 32 bytes does not ensure this
unpredictability. Set
\[
  (s_0,c_0)=
  \operatorname{split}_{32}\!\left(
    \mathrm{BLAKE2b}_{512}
      (\texttt{ZcashSA\_Issue\_V1},S)\right).
\]
The first half is a 32-byte secret and the second is a chain code.
For a hardened index \(i\ge 2^{31}\), define
\[
  (s',c')=\operatorname{split}_{32}\!\left(
     \mathsf{PRFexpand}_c
       (\mathtt{0x81}\,\|\,s\,\|\,\operatorname{enc}_4(i))\right).
\]
The byte \(\mathtt{0x81}\) is hexadecimal; the empty optional
lead and tag arguments of ZIP 32 contribute no bytes here.
Apply this step along
\[
  m_{\mathsf{Issuance}}/227'/\mathsf{coin\_type}'/0'.
\]
On mainnet the coin-type index is \(133'\).
The account index is fixed to zero in this selected design.
The final 32-byte secret is \(\mathsf{isk}\).

This is derivation of a \emph{byte string}. Interpreting that string
as the secp256k1 secret scalar below uses big-endian order. It is not
the Pallas little-endian scalar conversion used by RedPallas.
If the result is not a valid secp256k1 secret, key derivation fails
at the public-key step. ZIP 227 calls for discarding the key and
using a different issuance secret; it does not specify an automatic
hierarchical-key-derivation retry convention preserving the fixed path.

\subsection{The secp256k1 instance}

Issuance uses the Schnorr construction from \emph{Crypto Guide},
\S\emph{Knowledge, simulation, and Schnorr}, instantiated by
\href{https://github.com/bitcoin/bips/blob/master/bip-0340.mediawiki}
{BIP 340}. This is a different curve from Pallas.

The curve is
\[
  y^2=x^3+7\pmod {p_s},\qquad p_s=2^{256}-2^{32}-977,
\]
and its group order is the integer whose hexadecimal encoding is
\[
\begin{split}
 n_s={}&\mathtt{FFFFFFFFFFFFFFFFFFFFFFFFFFFFFFFE}\\[-1mm]
       &\mathtt{BAAEDCE6AF48A03BBFD25E8CD0364141}.
\end{split}
\]
The displayed lines concatenate; they are not an addition.
The standard generator \(G_s\) has coordinates
\[
\begin{split}
 x(G_s)={}&\mathtt{79BE667EF9DCBBAC55A06295CE870B0702}\\[-1mm]
          &\mathtt{9BFCDB2DCE28D959F2815B16F81798},\\
 y(G_s)={}&\mathtt{483ADA7726A3C4655DA4FBFC0E1108A8FD}\\[-1mm]
          &\mathtt{17B448A68554199C47D08FFB10D4B8}.
\end{split}
\]
Again concatenate each coordinate's lines. The arithmetic formulas
are those of \emph{Math Guide}, \S\emph{Pasta curves and their encodings},
with the curve constant changed from \(5\) to \(7\).
The curve's prime group order and discrete-logarithm hardness are
properties and assumptions of this selected standard curve, not
consequences of merely writing its equation.

Let \(\mathrm{BE}_{32}\) encode an integer in exactly 32 big-endian
bytes. A public key contains only the \(x\)-coordinate of its point.
To decode it, reject \(x\ge p_s\), compute
\[
  c=x^3+7\bmod p_s,\qquad y=c^{(p_s+1)/4}\bmod p_s,
\]
and reject unless \(y^2=c\). Of \(y\) and \(-y\bmod p_s\), select
the even representative. The exponent is a square-root candidate
because \(p_s\equiv3\pmod4\); the explicit square check rejects
non-residues. This operation is \(\mathsf{lift}_x\).

Given \(d'=\operatorname{int}_{\rm BE}(\mathsf{isk})\), reject unless
\(1\le d'<n_s\). The public key is
\(\mathsf{ik}=\mathrm{BE}_{32}(x([d']G_s))\).
For signing, replace \(d'\) by \(n_s-d'\) if its public point has
odd \(y\). The resulting \(d\) satisfies
\([d]G_s=\mathsf{lift}_x(\mathsf{ik})\).

\subsection{Tagged hashing and a complete issuance signature}

\emph{Crypto Guide}, \S\emph{Hash functions and key derivation}, constructs
SHA-256. BIP 340 gives each use a tagged wrapper:
\[
  H_{\mathsf{tag}}(m)=
  \mathrm{SHA256}\bigl(
    \mathrm{SHA256}(\mathsf{tag})\,\|\,
    \mathrm{SHA256}(\mathsf{tag})\,\|\,m\bigr),
\]
where the tag is its literal UTF-8 byte sequence.
For the issuance transaction's 32-byte signature digest \(M\),
the selected signing algorithm is as follows.
\begin{enumerate}
  \item Obtain the normalised scalar \(d\) and public key \(\mathsf{ik}\).
    Set
    \[
      t=\mathrm{BE}_{32}(d)\mathbin{\mathsf{xor}}
          H_{\texttt{BIP0340/aux}}(0^{32}).
    \]
    Here \(0^{32}\) is 32 zero bytes, as required by ZIP 227.
  \item Set
    \[
      k'=\operatorname{int}_{\rm BE}
        \bigl(H_{\texttt{BIP0340/nonce}}
           (t\,\|\,\mathsf{ik}\,\|\,M)\bigr)\bmod n_s.
    \]
    Reject \(k'=0\). If \([k']G_s\) has odd \(y\), replace \(k'\)
    by \(n_s-k'\). Call the result \(k\), and let \(R=[k]G_s\).
  \item Compute
    \[
      e=\operatorname{int}_{\rm BE}
       \bigl(H_{\texttt{BIP0340/challenge}}
        (\mathrm{BE}_{32}(x(R))\,\|\,\mathsf{ik}\,\|\,M)\bigr)
           \bmod n_s,
      \qquad s=k+ed\bmod n_s.
    \]
    Form \(\sigma=\mathrm{BE}_{32}(x(R))\,\|\,\mathrm{BE}_{32}(s)\).
    Verify it with the following algorithm before returning it;
    abort if this self-check fails.
\end{enumerate}

Verification first checks the public-key and signature lengths.
Lift the public key to the even-\(y\) point \(P\). Parse the signature
as integers \(r,s\), and reject \(r\ge p_s\) or \(s\ge n_s\).
Recompute \(e\) from \(\mathrm{BE}_{32}(r)\), the public key, and \(M\).
Set \(R'= [s]G_s-[e]P\). Accept precisely when \(R'\) is nonidentity,
has even \(y\), and has \(x(R')=r\).

Correctness is the identity
\[
  [k+ed]G_s-[e]([d]G_s)=[k]G_s.
\]
The parity normalisations ensure that the verifier uses the same
representatives as the signer. Security additionally requires
the computational and hash assumptions of Schnorr; correctness
alone does not imply unforgeability. Fixed auxiliary bytes do not
make the nonce public: its derivation includes the secret \(d\).
They do remove the additional fault and side-channel protection that
fresh auxiliary randomness can provide.

The scheme signature \(\sigma\) is 64 bytes. At ZIP 227's scheme-tagged
interface,
\[
  \mathsf{issuer}=\mathtt{0x00}\,\|\,\mathsf{ik},\qquad
  \mathsf{issueAuthSig}=\mathtt{0x00}\,\|\,\sigma,
\]
which are respectively 33 and 65 bytes. This does not specify the
outer transaction field framing; that is part of the carrier
boundary discussed later.

\subsection{From a description to an asset base}

An asset description \(D\) is a nonempty byte sequence, preferably
well-formed UTF-8. Consensus stores its digest, not a requirement to
publish a particular human-readable description.
Define
\[
\begin{aligned}
 h_D &=
  \mathrm{BLAKE2b}_{256}(\texttt{ZSA-AssetDescCRH},D),\\
 \mathsf{AssetId} &= (\mathsf{issuer},h_D),\\
 B_{\rm id} &= \mathtt{0x00}\,\|\,\mathsf{issuer}\,\|\,h_D,\\
 h_A &= \mathrm{BLAKE2b}_{512}(\texttt{ZSA-Asset-Digest},B_{\rm id}),\\
 A &= \GroupHash^{\mathbb P}
       (\texttt{z.cash:OrchardZSA},h_A).
\end{aligned}
\]
The BLAKE2b personalisation strings are zero-padded to their
16-byte parameter fields, as in \emph{Crypto Guide},
\S\emph{Hash functions and key derivation}, which also constructs
the exact Pallas \(\GroupHash\).

The first byte of \(B_{\rm id}\) is an asset-identifier version.
The next byte is the signature-scheme byte inside the issuer.
They happen both to be zero, but perform different jobs.
The total length is \(1+33+32=66\) bytes. \(h_A\) has 64 bytes,
and \(A^*\) has 32 bytes.

An asset name is not an identity. Two issuers may hash the same text
and still obtain different identifiers because their issuer fields
differ. Conversely, two bytewise different descriptions chosen to
collide in \(h_D\) would define the same identifier under one issuer.
Avoiding that attack is a collision-resistance assumption, not a
type-system guarantee. Wallets must identify assets by authenticated
identifiers or user-established names bound to them, not by an
untrusted description alone.

\subsection{The algebraic requirement on asset generation}

All nonidentity Pallas points generate the same cyclic group.
There is therefore some scalar \(a\) such that \(A=[a]V\), and
some scalar relation between any two asset bases. The construction
does not make those relations cease to exist. It must make useful
nontrivial relations computationally infeasible to find.

In particular, issuers must not choose a scalar \(a\) themselves
and declare \([a]V\) to be their asset base. Their choice is an
identifier preimage; the independently domain-separated group hash
determines the base. The later conservation argument assumes that
an efficient adversary cannot find a nonzero relation among these
derived bases and \(R_v\), within its permitted choices and queries.
Modelling the hash-to-group outputs as independent random group
points gives the usual discrete-logarithm-based justification.
The concrete hash is an instantiation of that model, not a proof
of its own hardness.

The map from 512-bit digests to a roughly 255-bit group cannot be
injective. Nor is the group hash's output type, by itself, a proof
that it avoids the identity or the native base. The note type
requires a nonidentity point, and native \(V\) has its own branch.
A complete deployment must resolve exceptional generated points
and asset-identity collisions consistently. It must not treat
a negligible-probability edge case as an unspecified parser choice.

\section{Issuance state and issue notes}
\label{sec:zsa-issuance}

\subsection{The public creation record}

An issue note contains
\[
  (d,\pkd,v,A,\rho,\mathsf{rseed}).
\]
Here \(d\) is an 11-byte diversifier, \(\pkd\) a nonidentity
diversified transmission key, \(v\in[0,2^{64}-1]\) an integer amount,
\(A\) the asset base, \(\rho\in\F_p\), and \(\mathsf{rseed}\) 32 bytes
of note randomness. The recipient construction is inherited from
\emph{Ironwood Guide}, \S\emph{Keys and viewing capabilities}, and \emph{Wallet Guide},
\S\emph{Addresses and receiver selection}.

Unlike an encrypted transfer output, these issue-note fields are
available to validators and chain observers. Anyone can reconstruct
its commitment once the note-format derivations are instantiated.
The issuer cannot conceal a larger minted value inside the commitment
while publicly accounting for a smaller one: the validator recomputes
the commitment from the public amount.

The recipient's spending authority is still private for an ordinary
issue note. Knowing its commitment opening is not knowing its
nullifier key or spending key. Later unlinkability nevertheless
depends on the inherited nullifier, proof, and transaction-privacy
assumptions; it is not an unconditional consequence of putting
the note into a common tree.

\subsection{Freshness derived from the transaction}

An \emph{Action Group} is an ordered nonempty list of transfer Actions
sharing a public anchor, spend/output/custom-asset enable flags, an
expiry bound, a public burn list, and a proof covering the group;
each Action also has its spend-authorisation signature. This is the
logical grouping required by the selected design's referenced
carrier, not a claim that the withdrawn carrier supplies a deployable
encoding. The selected issuance rules allow at most one such group.

Every issuance bundle accompanies a transfer bundle with an Action
from which a nullifier can be taken. Let \(n_0\) be the nullifier
of the first Action in the first Action Group. For issuance Action
index \(i\) and note index \(j\), both zero-based 32-bit integers,
the issued note uses
\[
 \rho_{i,j}=
 \operatorname{ToBase}\!\left(
   \mathsf{PRFexpand}_{\operatorname{enc}_{32}(n_0)}
    (\mathtt{0x84}\,\|\,\operatorname{enc}_4(i)
                      \,\|\,\operatorname{enc}_4(j))
 \right).
\]
Here \(\operatorname{ToBase}\) interprets the 64 output bytes
little-endian and reduces modulo \(p\). The first byte is \(0x84\).

This gives each note a transaction-dependent origin, rather than
letting an issuer reuse an arbitrary \(\rho\) in several transactions.
The input-nullifier uniqueness rules and the transaction signature
bind the construction to the containing transaction. The PRF's
key is public in this use; secrecy is not its purpose here.
Collision avoidance for derived \(\rho\) values is a computational
hash requirement, not literal injectivity of reduction modulo \(p\).

An \(\mathsf{IssueAction}\) contains a description hash \(h_D\),
a finalisation flag, and an ordered list of issue notes.
An issuance bundle contains one issuer identifier, an ordered list
of those Actions, and the issuer's signature over the \emph{whole
transaction} signature digest. It is not merely a signature over
the issuance amounts.

\subsection{The reference note}

For the first issuance of an asset, the first note must have value
zero and the default external diversified address derived from the
all-zero 32-byte Orchard spending key, at diversifier index zero.
Its address bytes are the following concatenation:
\[
\begin{array}{l}
\mathtt{cc36601959213b6b0cdb96a75c17c3a668a97f0d6a}\\
\mathtt{8c5ce164a518ea9ba9a50ea75191fd861b0ff10e62b0}.
\end{array}
\]
These 43 bytes are \(d\,\|\,\pkd^*\). The derivation, rather than
a visual resemblance to a human-readable address, gives the meaning
of this constant.

The reference is a real committed leaf but holds no value.
Its note data and its deliberately public spending key are known.
After its inclusion and an accepted anchor are available, anyone
can use the corresponding witness material to demonstrate the
existence of the asset in the split construction below. This
explains how that construction can obtain a publicly usable
membership witness without publishing a value-bearing spend key.

The reference's publicly known keys make it an exception to the
ordinary issue note's private-key argument. It must never be used
as a recipient for funds that are meant to be owned exclusively
by one person. Copying its opening does not mint units: a split
input contributes zero to the value equation.

\subsection{The state transition}

For a missing asset entry, read
\[
  I[A]=(0,\mathsf{false},\bot),\qquad
  \mathsf{MAX\_ISSUE}=2^{64}-1.
\]
The balance records \emph{issued minus burnt} units. It is a cap
on circulating ledger supply, not on the sum of all issuance events
over the asset's lifetime.

Start transaction validation with a working copy \(J\leftarrow I\).
The burn set contains distinct pairs \((A,b_A)\), with
\[
 A\ne V,\quad A\ne\mathcal O,\qquad
 1\le b_A\le 2^{63}-1.
\]
Every point must decode canonically. There may not be two burn
entries for the same base.

Process all burns before any issuance:
\begin{enumerate}
  \item For each \((A,b_A)\), require \(J[A].\mathsf{balance}\ge b_A\).
  \item Replace that balance by
    \(J[A].\mathsf{balance}-b_A\).
\end{enumerate}
The transaction's proof and binding signature, not this public
subtraction alone, establish that the sender actually supplies the
burnt units.

If there is an issuance bundle, require the accompanying Action
Group and a valid issuer signature over the transaction digest.
Process the issuance Actions in order:
\begin{enumerate}
  \item Derive \(A\) from the bundle's issuer and the Action's
    \(h_D\). Require \(J[A].\mathsf{final}=\mathsf{false}\).
  \item If \(J[A].\mathsf{reference}=\bot\), require a first note
    and check that it is the zero-valued reference note just
    specified. Store it as \(J[A].\mathsf{reference}\).
  \item For each note in order, validate its recipient and field
    encodings, compute its required \(\rho_{i,j}\), and check its
    equality with the supplied \(\rho\). Construct its note
    commitment from its public fields and the selected note
    derivations, rejecting any undefined commitment operation.
  \item Before adding \(v\), require
    \(J[A].\mathsf{balance}+v\le\mathsf{MAX\_ISSUE}\) as an
    integer inequality. Add \(v\) with checked arithmetic.
  \item If the Action finalises issuance, set
    \(J[A].\mathsf{final}\leftarrow\mathsf{true}\).
\end{enumerate}
Two issuance Actions for one asset are not made independent by
sharing a bundle. The second observes the first's state changes.
If the first finalises, a later Action for that asset fails.
An empty list cannot supply the mandatory first reference for a
new asset.

Only if the entire transaction passes all its proof, signature,
encoding, nullifier, anchor, and state checks does \(J\) become the
new issuance state. At activation the map is empty. Within a block,
each transaction sees the previous transaction's output map;
the first sees the parent's final map.

\subsection{Supply and finalisation invariants}

For one asset and any accepted branch prefix, let \(M_A\) be
the sum of public issue-note values and \(B_A\) the sum of public
burn values in that prefix. By induction on the ordered operations,
\[
  I[A].\mathsf{balance}=M_A-B_A,\qquad
  0\le M_A-B_A\le 2^{64}-1.
\]
The base case is the empty map. A burn subtracts the same amount
from both expressions, with its lower-bound check preserving
nonnegativity. An issuance adds the same amount, with its checked
upper bound preserving the cap. Transfer-only transactions alter
neither side. This proves the public accounting invariant exactly.

It does \emph{not} yet prove that privately spendable notes have
that total value. That stronger statement also requires the
transfer conservation argument, proof knowledge soundness, note
binding, and the nullifier uniqueness rules.

For example, issue \(1000\) units, burn \(100\), then issue \(50\).
The state balances are \(1000,900,950\), while lifetime issuance
is \(1050\). At a full supply cap, burning \(b\) and then issuing
\(b\) in one transaction is compatible with the cap, provided the
asset has not been finalised. Reversing the specified operation
order could change acceptance at the cap.

Finalisation is a monotone Boolean along a branch. Once true,
no later issuance Action for that asset is permitted. Burning
remains possible and does not clear the flag. Thus finalisation
limits future creation, not holders' transfers, and the outstanding
supply can still decrease.

\subsection{Appending the commitments}

The new tree leaves are ordered independently of the order in which
public supply changes are evaluated:
\begin{enumerate}
  \item Append transfer-output commitments in Action Group order
    and Action order.
  \item Append issue-note commitments in issuance Action order
    and note order.
\end{enumerate}
The reference note is included in the second part like any other
issue note. Commitment positions, paths, and anchors must agree
with this order. A public issue-note description is not a substitute
for a Merkle path against a recognised anchor.

\section{Asset-bearing notes and conservation}
\label{sec:zsa-conservation}

\subsection{Binding the asset into the note}

The conceptual note relation now opens
\[
  (d,\pkd,v,A,\rho,\psi,\rcm).
\]
For the native base \(A=V\), its commitment is the inherited Orchard
commitment, without an appended asset field.

For a custom asset define the bit string
\[
\begin{split}
  m_A={}&\gd^*\,\|\,\pkd^*\,\|\,\mathrm{LE}_{64}(v)
        \,\|\,\mathrm{LE}_{255}(\rho)\\
       &\,\|\,\mathrm{LE}_{255}(\psi)\,\|\,A^*,
\end{split}
\]
where point encodings in this expression are their 256-bit forms.
Then
\[
\begin{split}
 \cmsym_A={}&
   \mathsf{SinsemillaHashToPoint}
      (\texttt{z.cash:ZSA-NoteCommit-M},m_A)
       +[\rcm]R_c,\\
 R_c={}&
   \GroupHash^{\mathbb P}
     (\texttt{z.cash:Orchard-NoteCommit-r},\epsilon).
\end{split}
\]
The empty string is \(\epsilon\). \emph{Crypto Guide},
\S\emph{Groups, commitments, and Sinsemilla}, gives the bit
chunking, incomplete-operation rejection, and the distinction
between the full point and its extracted \(x\)-coordinate.
The leaf is the inherited extracted commitment \(\cmx\).

For a custom note, the input has
\(256+256+64+255+255+256=1342\) bits before Sinsemilla padding.
Changing the asset changes committed input bits as well as the
hash domain relative to a native note. The randomness base \(R_c\)
is shared; the native branch remains byte-for-byte the native
commitment construction.

The proof witnesses \(A\) once and uses it for the old note,
the new note, and the value commitment. Three unconstrained
copies merely bearing the same name would not enforce this
requirement.

\subsection{Randomness and the remaining note-format parameter}

The note's \(\psi\) is derived by the inherited ZIP 2005 function
\[
  \psi=\operatorname{ToBase}\!\left(
    \mathsf{PRFexpand}_{\mathsf{rseed}}
       (\mathtt{0x09}\,\|\,\operatorname{enc}_{32}(\rho))\right).
\]
The split-only randomness below deliberately uses \(0x0A\), not
this \(0x09\) domain.

The selected ZIP pair calls the ZIP 2005 note-randomness derivation
with its placeholder ZSA lead byte. The core rule distinguishes
legacy \(0x02\) randomness from recoverable \(0x03\) randomness,
but neither fact assigns a ZSA format byte. In particular, one
cannot replace the placeholder by \(0x03\) just because it once
appeared in a withdrawn carrier.

Accordingly, the relation here takes a valid note opening
\((\rho,\psi,\rcm)\) as witness data. Deriving \(\rcm\) from the
public \(\mathsf{rseed}\) of an \emph{issue note}, or constructing
an interoperable encrypted transfer plaintext, additionally needs
the missing ZSA note-format instantiation. This is a concrete
boundary of the selected draft, not permission for validators to
choose their preferred derivation.

\subsection{A private multibase value commitment}

For an ordinary Action with old and new values \(v_o,v_n\), set
\[
  C=[v_o-v_n]A+[r]R_v.
\]
The amount difference is a signed integer, subsequently interpreted
modulo \(q\) for scalar multiplication. The circuit's range checks
must establish that the note values are actual 64-bit integers
before this scalar interpretation.

There is one randomness scalar \(r\) for each Action. If \(r\) is
uniform in \(\F_q\), then for any fixed \(A\) and value difference,
\(C\) is uniform in \(\mathbb P\): translation by the fixed amount
term permutes the group. Thus this one commitment perfectly hides
its amount term under an independently uniform mask.

That marginal statement is not a proof of joint transaction privacy.
The transaction reveals a linear combination of commitments through
its binding key, contains note commitments and ciphertexts, and may
publish issuance or burn information. Security for all those
correlated objects needs the inherited proof and encryption
properties and appropriately joint randomness assumptions.

\subsection{The transaction equation}

Let \(\Delta_i=v'_{o,i}-v_{n,i}\), where \(v'_{o,i}\) is the
effective input value: the real value for a normal input and
zero for a split input. Then
\[
  C_i=[\Delta_i]A_i+[r_i]R_v.
\]
Let \(b_{\rm nat}\) be the native value leaving this bundle,
including its contribution towards the transaction fee, and let
\((A,b_A)\) range over the public burn set. The binding public key is
\[
  K_{\rm bind}
    =\sum_i C_i-[b_{\rm nat}]V-\sum_A[b_A]A.
\]
For an honest balanced transaction,
\[
  K_{\rm bind}=[k_{\rm bind}]R_v,\qquad
  k_{\rm bind}=\sum_i r_i.
\]
The binding signature uses the inherited binding-signature generator
\(R_v\), not the spend-authorisation generator or an issuer key,
and signs the complete transaction signature digest.

Collecting terms by asset gives
\[
 K_{\rm bind}
    =[\textstyle\sum_i r_i]R_v+\sum_A[\delta_A]A,
\]
where native \(V\) is included among the \(A\)'s, and
\[
 \delta_A=
 \begin{cases}
   \sum_{i:A_i=V}\Delta_i-b_{\rm nat},&A=V,\\
   \sum_{i:A_i=A}\Delta_i-b_A,&A\ne V,
 \end{cases}
\]
with \(b_A=0\) when no burn entry exists.

The equality is algebra. The claim that a successful adversary
cannot leave some \(\delta_A\ne0\) is computational. If the
Action proofs permit extraction of their values, bases, and masks,
and the signature argument permits extraction of a scalar \(k\)
for \(K_{\rm bind}\), an imbalance gives the known relation
\[
   \sum_A[\delta_A]A+
       [\textstyle\sum_i r_i-k]R_v=\mathcal O.
\]
This contradicts the no-efficient-relation assumption when at
least one coefficient is nonzero modulo \(q\).

Binding alone is insufficient for this argument: a binding
commitment does not automatically reveal its opening to a
reduction. Nor does ordinary fixed-key signature unforgeability
by itself establish extraction for an adversarially constructed
binding key. The required joint knowledge and composition
argument is the same kind of obligation identified in \emph{Ironwood Guide},
\S\emph{Balance, digests, and authorisation}, now with multiple
asset bases. This volume derives the relation but does not claim
a new proof of the complete Halo/signature composition.

\subsection{From field balance back to integer balance}

The relation argument establishes \(\delta_A=0\bmod q\).
To conclude \(\delta_A=0\) as an integer, prove an absolute bound
\(|\delta_A|<q\).

If there are at most \(N\) Actions, each effective input and output
lies in \([0,2^{64}-1]\). Therefore
\[
  \left|\sum_{i:A_i=A}\Delta_i\right|
       \le N(2^{64}-1).
\]
For a custom asset, a sufficient bound is
\[
  N(2^{64}-1)+(2^{63}-1)<q.
\]
For the native asset use the permitted absolute native balance
instead of the custom burn maximum. For example, even the very
loose \(N\le2^{32}-1\) would make these quantities below \(2^{97}\),
far below the Pallas group order. This is a sufficient arithmetic
condition, not a claim that the unspecified carrier actually
allows billions of Actions.

A complete transaction format must supply an explicit count or
size bound from which this condition follows. Range checks on
individual notes do not, by themselves, prevent arbitrarily
many field-wrapped terms from summing to zero.

\subsection{A two-asset check}

Suppose an Action spends \(80\) units of asset \(A\) and creates
\(50\), a second contributes a zero effective input and creates
\(20\) units of \(A\), and the burn set contains \((A,10)\).
The asset equation is
\[
  (80-50)+(0-20)-10=0.
\]
A separate native Action may pay a fee. Its native balance does
not alter this equation.

Replacing the second output with \(20\) units of another asset
\(B\) leaves nonzero per-asset coefficients. Even if the ordinary
integers appear to add up, the group equation contains different
bases. Acceptance would require a useful relation between those
bases and the randomness base, unless a proof or range constraint
has already failed.

\section{Split inputs and the counterfeiting boundary}
\label{sec:zsa-split}

\subsection{Why ordinary zero-value padding is insufficient}

An Orchard Action has one input position and one output position.
To turn one asset note into two output notes, the transaction needs
a second input position. For the native asset, a zero-value dummy
input can omit Merkle membership: the native base is a fixed
protocol constant and the dummy contributes no money.

If that exception were extended to an arbitrary custom base,
the prover could introduce a point with a known relation to \(V\).
Consider the hypothetical invalid choice \(A_{\rm evil}=-V\).
Create one native output of value \(w\) and one custom output of
value \(w\), both with dummy zero inputs. Their amount terms are
\[
  [-w]V+[-w]A_{\rm evil}
     =[-w]V+[w]V=\mathcal O.
\]
The binding-key equation could then hold although the transaction
created native value from nothing. All values can be small and
in range. The error would be admitting an attacker-chosen base,
not overflowing a machine integer.

This is an attack on the \emph{hypothetical weakened relation},
not on the complete split construction. It shows why merely
asserting that each Action uses the same base on its two sides
is not enough.

\subsection{Authenticate the base without consuming value}

A split input carries the opening and membership path of an existing
custom-asset note, with a private Boolean \(s=1\).
The old commitment is opened with its \emph{actual} old value \(v_o\).
Only its contribution to the value commitment is zeroed:
\[
  v'_o=(1-s)v_o,\qquad
  C=[(1-s)v_o-v_n]A+[r]R_v.
\]
When \(s=0\), the ordinary value equation is recovered.
When \(s=1\), changing a large stored \(v_o\) cannot create spendable
input value because the multiplier is exactly zero.

Let \(\nu\) be the private native-asset indicator. The relation
enforces
\[
  \nu\in\{0,1\},\quad
  \nu=1\Longleftrightarrow A=V,\quad
  s\in\{0,1\},\quad s\nu=0.
\]
Thus split inputs are custom-asset inputs, not another native
dummy-note encoding.

Write \(\mathsf{PathOK}\) for the inherited permitted path predicate
from \emph{Ironwood Guide}, \S\emph{Actions and the proof relation}.
It permits noncanonical 255-bit intermediate representatives that
reduce to the assigned field values, and the zero-hash alternatives
arising from exceptional Sinsemilla addition. It is not merely a
canonical path recomputation and root comparison.
The old-note membership condition is
\[
  (v_o=0\ \land\ \nu=1)
       \quad\lor\quad
  \mathsf{PathOK}(\cmx_o,\mathsf{path},
                        \mathsf{position},\mathsf{anchor}).
\]
All custom inputs require membership, including a zero-valued
reference note and a split copy of a previously spent note.
Whether the old note's ordinary nullifier is already spent is
irrelevant to its use as a split witness; the split creates a
different nullifier and supplies no old value.

The inherited key-consistency and spend-authorisation obligations
must also be discharged. A Merkle path and a guessed opening of
someone else's private note are not sufficient. The selected
ZIP's informal freedom to reuse an old note does not remove those
inherited proof equations. One can use one's own available
witnesses, obtain the required cooperation, or use the deliberately
public reference-note keys. The last option provides a universal
usable source without assuming access to arbitrary holders' keys.

\subsection{A fresh nullifier for the non-spending copy}

A normal spend uses the note's fixed nullifier. Reusing that
nullifier for every split would either expose repeated use of
the same source or cause the spent-nullifier set to reject later
splits. The split therefore introduces independent randomness.

Sample a fresh 32-byte seed \(u\) and compute
\[
  \psi_{\rm split}=\operatorname{ToBase}\!\left(
     \mathsf{PRFexpand}_u
       (\mathtt{0x0A}\,\|\,\operatorname{enc}_{32}(\rho_o))\right).
\]
Let
\[
 L=\GroupHash^{\mathbb P}(\texttt{z.cash:Orchard},\texttt{L}),
\]
and let \(K\) be the inherited Orchard nullifier base. The relation
uses a witnessed \(\psi_{\rm nf}\) constrained by
\[
   (1-s)(\psi_{\rm nf}-\psi_o)=0.
\]
For ordinary inputs this forces the original note randomness;
for a split the honest prover uses \(\psi_{\rm split}\).
The nullifier equation is
\[
 \nfsym=
  \Extract\!\left(
    [\,(\mathsf{Poseidon}(\nksym,\rho_o)+\psi_{\rm nf})
                               \bmod p\,]K
      +\cmsym_o+[s]L\right).
\]
The sum inside the brackets is first computed in the Pallas base
field, then its canonical integer is used for scalar multiplication.
It is not initially an addition modulo the group order \(q\).
The Poseidon nullifier PRF and \(\Extract\) are the inherited
functions described in \emph{Ironwood Guide},
\S\emph{Nullifiers and authenticated existence}.

The old commitment still uses \(\psi_o\), not
\(\psi_{\rm split}\). Otherwise a prover would change the leaf
whose membership it had supposedly demonstrated.
For the new note, the inherited Action linkage uses
\(\rho_n=\nfsym\).

The \(L\) offset is a group-level domain separation. It does not
reserve a visibly separate range of nullifier bytes: extraction
can map distinct points to the same field value. Independence
and collision resistance here are computational properties,
not an assertion that normal and split nullifiers occupy disjoint
sets.

The seed-to-\(\psi_{\rm split}\) expansion is an honest-prover
randomness algorithm, not an additional BLAKE2b computation proved
inside the Action circuit. The circuit permits a witnessed split
\(\psi_{\rm nf}\) subject to the stated relation. Consensus still
checks every emitted nullifier, including split nullifiers, for
duplicates within the transaction and against the branch's
spent-nullifier set.

\subsection{Provenance is an inductive invariant}

Assume note-commitment binding across the native and custom
branches, key-commitment binding for the inherited ownership
relation, authenticated tree membership under the permitted
representations, and knowledge soundness of the Action proofs.
Also assume the Sinsemilla exceptional-case hardness from
\emph{Crypto Guide}, \S\emph{Groups, commitments, and Sinsemilla}:
efficiently causing the permitted failure output \(\bot\), rather
than an ordinary commitment or hash, yields an infeasible generator
relation. This assumption is needed both for note/key commitment
alternatives and for zero-hash path acceptance. A proof of the
weakened relation alone does not remove those alternatives.
Outside these explicitly bounded failure events, consider accepted
custom-asset notes in their order of creation.

An issue note's base is publicly derived from the issuer and
description hash and checked by the issuance validator.
Every custom transfer output shares its base with an input note
whose membership is proved. That input therefore supplies an
earlier committed occurrence of the base. By induction, the
base traces to a publicly validated issuance. A split input
preserves this argument because its membership condition is
unchanged.

The induction concerns \emph{asset origin}, not ownership of
all previous notes or preservation of their values. The latter
is discharged by the effective-value equation and the binding
signature. It also concerns the accepted branch and recognised
anchors: a path into an unrelated or rejected tree does not
establish provenance.

This separates two obligations that are easy to conflate.
Membership prevents introduction of an arbitrary relational base;
balance prevents creation of extra units of a legitimate base.
Neither condition alone implies the other.

\section{The Action relation and transaction validation}
\label{sec:zsa-validation}

\subsection{Public statement and private witnesses}

The public statement contains the inherited anchor, value
commitment, input nullifier, randomised spend-authorisation key,
output commitment, and spend/output enable flags, together with
the custom-asset enable flag \(e_A\).
The asset base and split/native indicators remain private.

The witness contains the old and new note openings, the old-note
path and position, the inherited ownership and key-randomisation
witnesses, the value-commitment mask, and the additional
\(A,\nu,s,\psi_{\rm nf}\). The new circuit must connect these
objects, rather than checking disconnected subrelations:
\begin{enumerate}
  \item Validate each point and scalar in its proper type, and
    constrain both note values to 64-bit integers.
  \item Constrain the inherited diversified-key, incoming-viewing-key,
    and randomised spend-authorisation-key equations. A split
    is not a blanket exception to key consistency.
  \item Open the old and new note commitments using the one
    shared asset base and the appropriate native/custom branch,
    with the inherited exceptional Sinsemilla alternatives.
  \item Enforce \(\nu=1\) exactly when \(A=V\), require Boolean
    \(s,\nu\), and require \(s\nu=0\).
  \item Require \(\mathsf{PathOK}\) unless it is a native zero-value
    dummy. Check the conditional normal/split nullifier and the
    new-note \(\rho\) linkage.
  \item Enforce the value commitment with effective old value
    \((1-s)v_o\).
  \item Enforce the inherited spend/output enable conditions and
    require \(e_A=0\Rightarrow\nu=1\).
\end{enumerate}
The last rule prevents a prover from hiding custom-asset use in
a proof whose public enable flag disables it. Selecting the flag
for a transaction or upgrade is a consensus policy obligation,
not a choice that can safely be delegated to an unconstrained
private witness.

\subsection{How the Boolean conditions become field equations}

\emph{Halo 2 Guide}, \S\emph{From a relation to a circuit}, explains selectors,
constraints, and range checks. Here a Boolean \(s\) is enforced
by \(s(s-1)=0\) in the prime field. Similarly enforce
\(\nu(\nu-1)=0\).

Let \(a_x,a_y\) be the canonical affine coordinates of the
nonidentity \(A\), and \(v_x,v_y\) those of \(V\).
For the native direction use
\[
   \nu(a_x-v_x)=0,\qquad \nu(a_y-v_y)=0.
\]
For the non-native direction, witness \(u_x,u_y\) and enforce
\[
 (1-\nu)\bigl((a_x-v_x)u_x-1\bigr)
          \bigl((a_y-v_y)u_y-1\bigr)=0.
\]
When \(\nu=0\), at least one coordinate difference must have
an inverse. If both are zero, both parenthesised factors are
\(-1\), so the equation fails. If one is nonzero, its inverse
makes the product zero. This enforces inequality of points,
not just inequality of their \(x\)-coordinates.

In the ordinary nonexceptional path calculation, let
\(r_{\rm path}\) be the computed root. The membership gate is
\[
   (v_o+1-\nu)(r_{\rm path}-\mathsf{anchor})=0.
\]
Because \(v_o\) is range-constrained and \(2^{64}<p\), the first
factor is zero exactly when \(v_o=0,\nu=1\). Without the range
check it could also vanish through an unintended field value.
This gate explains the native-dummy bypass; the complete inherited
\(\mathsf{PathOK}\) predicate retains the representation and zero-hash
allowances stated above. It is not strengthened to canonical root
equality by this illustrative equation.

The inherited enable conditions are
\[
   v_o(1-e_{\rm spend})=0,\qquad
   v_n(1-e_{\rm output})=0,
\]
and the new condition is
\[
   (1-e_A)(1-\nu)=0.
\]
Notice that the spend-enable equation uses the \emph{old note
value}, not its zeroed split contribution. Copying a nonzero
old note still requires spends enabled under the inherited
condition. A zero-valued public reference avoids that particular
nonzero-value requirement.

These equations explain the new gates. They do not replace the
elliptic-curve, Sinsemilla, Merkle, and key-consistency gadgets
they connect, or prove that a particular compiled circuit has
wired them all correctly.

\subsection{Complete validation order at the defined interface}

A validator must combine the following checks, using the generic
atomic state-transition contract in Consensus and the inherited
payment checks in \emph{Ironwood Guide}, \S\emph{Pool rules and complete validation}.
\begin{enumerate}
  \item Select the rules for the containing branch and activation
    context. Parse canonical transaction and field encodings,
    reject inconsistent counts, and enforce resource bounds.
    Under the selected ZIP 227 rule, the Action Group count is
    zero or one and its group expiry height is zero.
  \item Check the bundle flags, recognised anchor, and every
    nullifier's uniqueness within the transaction and against
    the branch state. A split nullifier is not exempt.
  \item Reconstruct the exact public proof instances from the
    transaction fields and verify the proof under the matching
    circuit/verifying key.
  \item Compute the transaction signature digest and verify
    every spend-authorisation signature under its randomised key.
    Recompute the multibase binding key and verify its binding
    signature.
  \item Check the distinct public burn entries and run the
    burn-first issuance-state transition. If issuance is present,
    verify its issuer signature, reconstruct every issue note,
    and check the reference, supply, and finalisation rules.
  \item Check the inherited native transaction balance and fee
    conditions, as well as all remaining contextual transaction
    rules.
  \item Commit all state changes together: spent nullifiers,
    transfer and issuance commitment leaves in the specified order,
    and the new issuance map. Otherwise change none of them.
\end{enumerate}
The expensive cryptographic checks can be batched or reordered
operationally when that preserves acceptance semantics. The
state effects cannot be partially applied merely because an
earlier check succeeded.

\subsection{What is still needed for a deployed endpoint}

Three unresolved interfaces in the selected pin prevent the above
from being a complete network implementation.
\begin{enumerate}
  \item \emph{Carrier and digests.} A non-withdrawn transaction
    encoding must frame all transfer, issuance, and burn fields,
    fix count/size limits, assign branch/version identifiers, and
    specify which exact bytes every signature and authorising
    commitment covers.
  \item \emph{Note-format instantiation.} The ZSA plaintext lead
    byte, note-randomness derivation, encrypted plaintext lengths,
    and incoming/outgoing decryption behaviour must agree.
    In particular, validators must deterministically reconstruct
    issue-note commitments from the public fields.
  \item \emph{Circuit and activation.} The complete relation must
    be implemented and its verifying key selected by explicit
    consensus rules, including the meaning of the enable flag
    and exceptional invalid asset-generation outcomes.
\end{enumerate}
The new group equations do not fill these gaps automatically.
Nor do they supply a new proof of the complete multibase
knowledge-soundness and authorisation composition.

Within the specified interface, the critical path is now explicit:
issuer identity determines a hashed asset base; public issuance
places that base in authenticated notes; every custom transfer
preserves such provenance; constrained value commitments and
authorisation enforce conservation; public burns and issuance
update a branch-local supply map. Every step has a distinct
cryptographic or state-machine obligation, and no step can be
replaced by trusting a displayed asset name.

\end{document}
